%% sections/03-mechanics.tex

\section{The Two Resource Mechanics}
\label{sec:mechanics}

\subsection{Campaign 4: Supply Die + Between-Session HP Carry-Over}

\paragraph{Design.}
Campaign 4 (the Thorngate Watch siege scenario) implemented two interlocking resource
mechanics:

\begin{enumerate}
  \item \textbf{Between-session HP carry-over:} At the end of each session, HP is recorded
    and carried forward. At the start of the next session, each PC recovers 1 hit die of
    HP (rolled or average), capped at their maximum. No full long rest is available until
    the siege breaks (at campaign end). Spell slots recover normally (short rest = 1 warlock
    slot; one short rest per session allowed; no long rest slots until finale).

  \item \textbf{Supply Die:} The garrison's medical and supply resources are tracked as a
    single die that begins at d12 and degrades one category (d12$\to$d10$\to$d8$\to$d6$\to$d4)
    after each session in which a combat encounter occurs. At each die level, the garrison
    can provide a corresponding level of healing to the party.
\end{enumerate}

The two mechanics interact: HP carry-over creates cumulative pressure on character HP
decisions (when to use Mending Touch in session 3 affects how much HP is available in
session 6); Supply Die exhaustion removes external healing as a safety net. By Session 6,
the party is fighting at reduced HP with zero garrison healing support.

\paragraph{Design intent.}
The mechanics were designed to make every session's resource decisions visible across
the full campaign: spell slots spent in Session 2 are slots not available in Session 7;
HP taken in Session 3 is HP that must be recovered (partially) across multiple sessions.
The design inverts the typical TRPG session structure, in which long rests reset resources
between sessions.

\paragraph{Emotional architecture.}
The Supply Die represents the garrison's wounded --- the people who need healing to
survive the siege. When the die is exhausted in Session 6, it is not just the party
that has run out of resources; it is the garrison that has no more medical supplies.
The resource depletion is felt as a narrative statement (``we are at our limit'') as
well as a tactical statement (``no healing available'').

\subsection{Campaign 6: Unit Cohesion Die}

\paragraph{Design.}
Campaign 6 (the Border Watch military unit) used a Unit Cohesion Die representing the
number of soldiers available to the party's unit:

\begin{itemize}
  \item Begins at d12 (full unit of 12+ soldiers).
  \item Degrades one category when unit members are killed or captured in combat.
  \item Degrades one category when a PC defies direct orders without unit consensus.
  \item Does not recover between sessions (unit members who are gone are gone).
\end{itemize}

When the die is rolled for morale checks or tactical coordination, the result is
compared to the number of soldiers available. As the die degrades, the unit's capacity
for coordinated action diminishes.

\paragraph{Distinction from Supply Die.}
The key distinction is what the resource represents. The Supply Die tracks materials
(healing supplies); the Cohesion Die tracks people (unit members). Material depletion
is felt as scarcity; people depletion is felt as loss. The party's relationship to
the Cohesion Die is characterized by active preservation: the party in Campaign 6
consistently made decisions that prioritized keeping unit members alive over tactical
advantage. This preservation behavior is itself a form of character-driven agency.

\paragraph{Emotional architecture.}
Each Cohesion Die degradation represents named unit members. In Campaign 6, each
degradation event was recorded with the name(s) of who was lost. By Session 7, the
die at d8 represented a unit from which 4 members had been lost. The party could
track their losses as a named list. This made the resource emotionally legible in a
way the Supply Die is not: it is harder to care about healing supplies than about
people you have worked with for six sessions.

\subsection{Mechanic Comparison: Design Matrix}

Table~\ref{tab:mechanic-comparison} compares the two resource mechanics on the design
properties proposed in our amplification hypothesis.

\begin{table}[t]
\caption{Supply Die vs.\ Unit Cohesion Die: Design Property Comparison}
\label{tab:mechanic-comparison}
\begin{tabular}{@{}lll@{}}
\toprule
Property & Supply Die (C4) & Cohesion Die (C6) \\
\midrule
Resource represents & Materials (healing supplies) & People (unit members) \\
Degradation trigger & Session with combat & Unit loss OR order-defiance \\
Legibility & Die-size visible (d12→exhausted) & Die-size visible (d12→d8) \\
Named loss & No (generic ``supplies'') & Yes (named unit members) \\
Irreversibility & Full (supply not recoverable) & Full (members not recoverable) \\
Preservation behavior observed & Minimal (use when needed) & Strong (active conservation) \\
Emotional architecture strength & Moderate & High \\
Combat pressure produced & High (HP deplete) & Moderate (unit attrition rate) \\
Character Agency score S5-7 & 10/10/10 & 10/10/10 \\
Panel CA quality notes & ``Character-inflected tactical'' & ``Character-primary'' \\
\bottomrule
\end{tabular}
\end{table}

Both mechanics are legible and irreversible. The critical distinction is in the
``named loss'' and ``preservation behavior'' rows. The Cohesion Die tracks named
individuals, which activates preservation behavior absent from the Supply Die campaign.
This preservation behavior is itself character-driven agency: it is the party choosing
to sacrifice tactical efficiency to protect specific named people.

In Campaign 6 sessions 4--6, the Tactician panel lens noted on three separate occasions
that the unit's defensive formation (forgoing offensive opportunities to hold unit
members in protected positions) was sub-optimal from a damage-per-round standpoint but
was explained in the session log as a character commitment rather than tactical
calculation: ``Varet extended the left flank at the cost of a flanking bonus because
the three soldiers holding that position included Carra, who had held the ridge all
night. She did not explain this.'' This is act-without-announcement applied to resource
management: the Cohesion Die made the party's commitment to their people visible as a
resource decision, not just as a narrative gesture.
